\documentclass[a4paper,12pt]{article}
\usepackage[a4paper,top=1.3cm,bottom=2cm,left=1.5cm,right=1.5cm,marginparwidth=0.75cm]{geometry}
\usepackage{setspace}
\usepackage{cmap}
\usepackage{mathtext}
\usepackage[T2A]{fontenc}
\usepackage[utf8]{inputenc}
\usepackage[english,russian]{babel}
\usepackage{multirow}
\usepackage{graphicx}
\usepackage{wrapfig}
\usepackage{tabularx}
\usepackage{float}
\usepackage{longtable}
\usepackage{hyperref}
\hypersetup{colorlinks=true,urlcolor=blue}
\usepackage[rgb]{xcolor}
\usepackage{amsmath,amsfonts,amssymb,amsthm,mathtools}
\usepackage{icomma}
%\mathtoolsset{showonlyrefs=true}
\usepackage{euscript}
\usepackage{mathrsfs}
\usepackage{float}
\setlength{\parindent}{1cm} % Устанавливает отступ в 1.5 см для всех абзацев
\usepackage{multicol}
\usepackage{caption}




\DeclareMathOperator{\sgn}{\mathop{sgn}}
\newcommand*{\hm}[1]{#1\nobreak\discretionary{}
	{\hbox{$\mathsurround=0pt #1$}}{}}

\title{\textbf{Отчёт о выполненой лабораторной работе \\ \text{Измерение вязкости воздуха по течению \\
в тонких трубках (1.3.3)}}}

\author{Дульмиев Муса Б01-507}
\date{март 2026}

\begin{document}
\maketitle


\begin{figure}[h!]
        \centering
        \includegraphics[scale=0.5]{expust.png}
        \caption{Экспериментальная установка}
 \end{figure} 

\begin{table}[!ht]
    \centering
    \begin{tabular}{|c|c|c|c|}
        \hline
        $T, K$ & $h, \text{мм}$ & $P, \text{Па}$ & $\sigma, \text{мН/м}$ \\ \hline
        $293,2 \pm 2$ & $167,3 \pm 1$ & $327$ & $86$ \\ \hline
        $298,0 \pm 2$ & $165,7 \pm 1$ & $323$ & $85$ \\ \hline
        $303,0 \pm 2$ & $165,0 \pm 1$ & $322$ & $85$ \\ \hline
        $308,3 \pm 2$ & $164,0 \pm 1$ & $320$ & $84$ \\ \hline
        $313,2 \pm 2$ & $164,7 \pm 1$ & $321$ & $84$ \\ \hline
        $318,1 \pm 2$ & $164,0 \pm 1$ & $320$ & $84$ \\ \hline
        $323,0 \pm 2$ & $162,3 \pm 1$ & $317$ & $83$ \\ \hline
        $327,9 \pm 2$ & $160,3 \pm 1$ & $313$ & $82$ \\ \hline
        $333,0 \pm 2$ & $159,5 \pm 1$ & $311$ & $82$ \\ \hline
    \end{tabular}
    \caption{Результаты измерений зависимости $\sigma(T)$}
    \label{tab:sigma_T}
\end{table}

\end{document}